
\sourcenote{Genre: admissible rule (tier 2). The family machinery is
stated in the \rulename{LiRPA} section's family template. This rule is
premise-free: $\Phi$ is arbitrary and never consulted.}

\begin{brecipe}[the \rulename{ReLU-Above} instance --- worked plumbing]
\label{rec:reluabove-instance}
The node:
\[
\inferrule*[Right=\rulename{ReLU-Above}]{
  \ReLU_{\{h\}\to\{y\}} \in H
}{
  \Phi \vdash y \geq h
}
\qquad\rightsquigarrow\qquad
\text{case analysis on }
(h \leq 0 \wedge y = 0) \vee (0 \leq h \wedge y = h).
\]
We walk the \emph{inactive} branch in full --- the family's honest
specimen of what ``one node unpacks to many'' means in catalogue
practice. Context $\Phi_{-}$ carries $h \leq 0$ and $y = 0$; the goal
is $y \geq h$.
\begin{enumerate}
\item $\Phi_{-} \vdash h \leq 0$ --- \rulename{Ax} + $\wedge$E.
\item $\Phi_{-} \vdash -h \geq 0$ --- the bound node's polarity-flip
  reading (\rulename{PBound}, reading (G)): the combination is the
  single component $h \leq 0$ with weight $1$, canonical polynomial
  $c = h$, read as $e \geq 0$ with $e = -h$.
\item $\Phi_{-} \vdash (-h) = (0 - h)$ --- \rulename{Norm}: equal
  monomial canonical forms over $\Q$, same wires, shared broadcast
  shape.
\item $\Phi_{-} \vdash (0 - h) \geq 0$ --- \rulename{EqElim} along
  step 3's equality (both sides shaped alike --- the transport gate is
  met), rewriting step 2's atom.
\item $\Phi_{-} \vdash 0 \geq h$ --- \rulename{ArithAtom} (fold): from
  $(t_1 - t_2) \bowtie 0$ to $t_1 \bowtie t_2$ at $t_1 = 0$, $t_2 =
  h$, $\bowtie = {\geq}$.
\item $\Phi_{-} \vdash 0 = y$ --- \rulename{EqSym} on the context's
  $y = 0$; then $\Phi_{-} \vdash y \geq h$ --- \rulename{EqElim} along
  $0 = y$ at $\Psi = (a \geq h)$.
\end{enumerate}
The \emph{active} branch (context $0 \leq h$, $y = h$) derives $h \geq
h$ by the same pattern --- \rulename{Norm} gives $(h - h) = 0$,
\rulename{PConstExact} the ground fact $0 \geq 0$, \rulename{EqElim} +
\rulename{ArithAtom} assemble $h \geq h$ --- and \rulename{EqElim}
along $h = y$ (wire-to-wire, one shared shape) closes the goal. Six
catalogue nodes to move one minus sign across a comparator: this is
why the family rules exist as single nodes, and why the wire format
prices them.
\end{brecipe}

\begin{bremark}[shape scope]
Steps 3 and 6 of the inactive branch (and the active branch's
\rulename{Norm} identity $(h - h) = 0$) transport along equalities
with a rank-$0$ constant side, so the displayed derivation
type-checks at rank-$0$ instances; the shape-general warrant is the
direct lemma, exactly as Remark~\ref{rem:relunn-shapescope} in the
\rulename{ReLU-NN} section states for the whole family.
\end{bremark}

\begin{bprop}[semantic warrant]\label{prop:reluabove-sound}
If the linkage holds, the sequent concluding $y \geq h$ is
semantically valid over any premise formula $\Phi$.
\end{bprop}
\begin{proof}
Fix an admissible valuation $\sigma$. The linkage extraction
(\texttt{reluLinked\_facts}) gives both tensors at one shared shape
and, per entry, either $p \leq 0 \wedge y = 0$ --- whence $y = 0 \geq
p$ --- or $y = p$ --- whence $y \geq p$ reflexively. The conclusion
atom compares the two tensors entrywise at their shared shape. That is
\texttt{sound\_reluAbove}, step for step.
\end{proof}

\begin{bremark}[the exact lower envelope, premise-free]
Together with \rulename{ReLU-NN} this rule states the exact pointwise
lower envelope of the activation, $y \geq \max(0, h)$, as two atoms a
producer can emit with \emph{no} box knowledge whatsoever. The
asymmetry with the upper side is structural, not an artefact: no
affine function dominates $\max(0, h)$ on all of $\Q$ (a candidate
with positive slope dives below $0$ as $h \to -\infty$; one with
slope $\leq 0$ falls below $h$ as $h \to +\infty$) --- so upper
bounds require a box and live in \rulename{LiRPA}.
\end{bremark}

\begin{anchors}
Recipe \ref{rec:reluabove-instance} warrant & \texttt{sound\_reluAbove}\\
family linkage extraction & \texttt{reluLinked\_facts} (see the \rulename{LiRPA} section)\\
polarity-flip reading (step 2) & \texttt{sound\_pBound} (the \rulename{PBound} section)\\
ground-fact leaf (active branch) & \texttt{sound\_pConstExact} (the \rulename{PConstExact} section)\\
payload structure & \texttt{PReLUAboveData} (two wires, linkage --- no premise field)\\
\end{anchors}
